\section{Mag (Magmas Finitos)}
\label{sec:mag}

\subsection{Definição da Categoria}

\subsubsection{Categoria $\mathbf{Mag}$}

\begin{definition}
    Os objetos de $\mathbf{Mag}$ são magmas finitos, ou seja, conjuntos finitos equipados com uma operação binária sem axiomas adicionais.
    Os morfismos são aplicações que preservam a estrutura do magma, isto é, para dois magmas $(M, \cdot)$ e $(N, \ast)$, um morfismo $f: M \to N$ satisfaz $f(a \cdot b) = f(a) \ast f(b)$ para todo $a, b \in M$.
\end{definition}

\begin{remark}
    Não são impostos axiomas adicionais na operação binária (associatividade, comutatividade, etc.).
\end{remark}

\subsection{Implementação Computacional}

\begin{lstlisting}[style=matlab]
% Definição de um magma
Mag = struct();
Mag.set = {1, 2, 3}; % Conjunto finito
Mag.op = @(a, b) mod(a + b, 3); % Operação binária

% Verificação de morfismo entre magmas
MorphismMag = @(M1, M2) @(f) ...
    all(arrayfun(@(a, b) M2.op(f(a), f(b)) == f(M1.op(a, b)), M1.set, M1.set));
\end{lstlisting}

\subsection{Transformações Naturais}
% Exemplo de lugar para implementação de transformações naturais
\begin{lstlisting}[style=matlab]
function ok = check_morphism(M1, M2, f)
    ok = all(arrayfun(@(a, b) M2.op(f(a), f(b)) == f(M1.op(a, b)), M1.set, M1.set));
end
\end{lstlisting}